\documentclass[10pt]{article}
\usepackage[margin=2cm]{geometry}
\usepackage[T1]{fontenc}
\usepackage{lmodern,microtype,booktabs,longtable,array,xcolor,amssymb,enumitem,hyperref}
\hypersetup{colorlinks=true,linkcolor=blue!50!black,urlcolor=blue!50!black}
\setlist{nosep,leftmargin=1.4em}
\newcommand{\decision}[1]{\par\medskip\noindent\fcolorbox{orange!70!black}{orange!6}{\parbox{\dimexpr\linewidth-2\fboxsep-2\fboxrule}{\textbf{Decision.} #1\par\smallskip
$\square$ Approve as is \qquad $\square$ Change (note below) \qquad $\square$ Discuss\par\smallskip Notes: \hrulefill\par\vspace{1.2em}\hrulefill}}\par\medskip}
\newcolumntype{P}[1]{>{\raggedright\arraybackslash}p{#1}}
\title{Exposome / EHR rebuild: what needs your review}
\author{Branch \texttt{feature/decision-model-rebuild}, jeonglab-bcm/exposome-ehr-review}
\date{Generated by \texttt{scripts/build\_review\_report.py}, 2026-10-05}
\begin{document}
\maketitle
\section*{Summary}
The repository was rebuilt on the POTS-phenotyping design: one broad PubMed query, facets as tags rather than
filters, and literature links found by decision models and confirmed by Claude with verbatim quotes.
The code is tested (132 offline tests) and the counts below are taken from the committed harvest. What only you
can decide are the four choices below. The model runs are built on Decisions~1 and~2, so changing those afterwards
means re-running; review them first.

\begin{center}\begin{tabular}{@{}clP{7.2cm}l@{}}\toprule
 & Decision & What to look at & Section\\\midrule
1 & Scope of the search & 4 query blocks; what is deliberately left out; 4 known misses & \ref{sec:scope}\\
2 & Seed papers & 8 papers the evidence graph is first built from & \ref{sec:seeds}\\
3 & Vocabulary & 67 terms in 6 groups; 9 flagged judgement calls & \ref{sec:vocab}\\
4 & Per-paper questions & 8 fields replacing the old checklist & \ref{sec:coding}\\\bottomrule
\end{tabular}\end{center}
Corpus: 15,913 articles, 6,951 tagged pediatric (prenatal, infant, child or adolescent). Section~\ref{sec:skip} lists what does not need your review.
\tableofcontents
\section{Decision 1: scope of the search}\label{sec:scope}
A paper enters the corpus if it matches any of the four blocks below (\texttt{config/query.yaml}). Nothing else is
filtered at search time: age, publication type and open access are recorded per paper and filtered later.

\begin{center}\begin{tabular}{@{}lP{8.6cm}rr@{}}\toprule
Block & What it retrieves & Papers & Only here\\\midrule
\texttt{exposome} & The exposome concept itself (Wild 2005) & 3,123 & 2,892\\
\texttt{ewas} & Agnostic, many-exposure association designs (Patel 2010) whose abstracts do not always say "exposome". & 297 & 113\\
\texttt{environment\_in\_health\_records} & The core of the old Tiers 2-4: an environmental or social exposure, studied in routinely collected individual-level data & 11,624 & 11,558\\
\texttt{vaccine\_as\_exposure} & Old Tier 5: a vaccine as the exposure and a health event as the outcome, in surveillance, registry, claims or cohort data & 1,116 & 1,109\\
\bottomrule\end{tabular}\end{center}

\paragraph{Deliberately left out.}
\begin{itemize}
\item \textbf{Occupational exposure.} PubMed's ``Environmental Exposure'' heading includes it; 1,647 papers came in
only that way, so the heading is used without its children (maternal, prenatal and inhalation exposure are kept).
\item \textbf{Research-cohort studies with no routinely collected data}, unless they say ``exposome'' or are
exposome-wide (e.g.\ a single chemical in one birth cohort). The old pipeline included some of these.
\item \textbf{Vaccine uptake, coverage and hesitancy} studies (vaccination as the outcome, not the exposure);
excluded by title only.
\end{itemize}

\paragraph{Known misses and the widenings declined.} Of the old pipeline's 30 EHR-positive papers, 11 are retrieved;
15 of the other 19 are out of scope (2 corrections, 10 programme reports, descriptive surveillance or attitude
surveys, 2 research cohorts, 1 economic estimate). Four are real misses:
\begin{itemize}
\item PMC4997457 (air pollution and asthma exacerbations, Quebec administrative data): the abstract names only
``hospitalization or emergency room visit''. Adding such phrases grows the corpus to about 19,800, mostly aggregate
time-series studies.
\item PMC8986305 (vaccine safety comparison): adding surveillance and pharmacovigilance terms to the vaccine block
recovers it but adds about 1,600 papers.
\item PMC4381500 (mortality after DTP/MMR eligibility) and PMC4113360 (prenatal bereavement, registry data).
\end{itemize}
\decision{Are these the right boundaries? In particular: include occupational exposure? include single-exposure
research cohorts? accept the emergency-visit widening (+3,900 papers) or the vaccine-surveillance widening
(+1,600)?}
\section{Decision 2: seed papers}\label{sec:seeds}
The first model run reads these papers paragraph by paragraph; their links start the evidence graph, as POTS's seven
seeds did. Only CC BY text may be sent to the hosted Jev model, so most seeds rely on the local models agreeing.

\begin{center}\begin{tabular}{@{}lrP{5.6cm}lr@{}}\toprule
Seed & PMID & Role & Licence & Paragraphs\\\midrule
wild2005 & 16103423 & Defines the exposome concept (commentary) & abstract only & 0\\
patel2010 & 20505766 & Agnostic EWAS method, NHANES, adults & CC BY & 49\\
maitre2018 & 30206078 & HELIX: pediatric multi-exposure birth cohort & CC BY-NC & 51\\
vrijheid2020 & 32579081 & Pediatric exposome-wide study of obesity & EHP & 58\\
brokamp2018 & 29126118 & Geocoding environmental exposures onto EHR & CC BY-NC & 32\\
hu2023 & 37089437 & EHR exposome-wide study (OneFlorida) & CC BY-NC & 28\\
correa2022 & 35970309 & Neighbourhood factors and childhood asthma & author manuscript & 33\\
macdonald2014 & 24914115 & Vaccine as exposure, linked administrative data & abstract only & 4\\
\bottomrule\end{tabular}\end{center}
Wild 2005 has no abstract in PubMed and no PMC full text, so it contributes no paragraphs; it stays a seed because
it defines the concept. The vocabulary's \texttt{source\_refs} were assigned from the seeds' abstracts, not yet
from reading them.
\decision{Keep these eight? Candidates to add: a pediatric EHR study of a chemical exposure (lead from EHR
screening results), a Vaccine Safety Datalink study with an open licence, and an adult EHR exposome-wide study
other than Hu 2023.}
\section{Decision 3: vocabulary}\label{sec:vocab}
Each term is a tag, not a category: a paper carries any number. ``Papers'' is how many corpus papers carry the tag
by either path (title/abstract query or NLM MeSH indexing). Rows marked \textbf{?} are judgement calls.
\subsection{Exposure}
An environmental, chemical, social or medical factor a person encounters, studied as a possible cause of a health outcome.

\begin{longtable}{@{}P{3.6cm}P{9.4cm}r@{}}\toprule Term & Definition (first sentence) & Papers\\\midrule\endhead
Particulate matter air pollution & Ambient fine or coarse particles (PM2.5, PM10, black carbon, ultrafine particles), usually assigned from a residential address & 1,318\\
Traffic-related air pollution & Nitrogen dioxide, nitrogen oxides and other markers of road traffic, including distance to a major road and traffic density & 738\\
Ozone & Ground-level ozone, a secondary photochemical air pollutant & 430\\
Wildfire smoke & Smoke from wildfires, bushfires or landscape fires, as an episodic air pollution exposure & 126\\
Indoor air and household fuel & Indoor sources: gas stoves, household solid fuel, mould and dampness, indoor allergens and radon & 276\\
Tobacco smoke & Active smoking, maternal smoking in pregnancy, or second-hand (environmental) tobacco smoke, including e-cigarette aerosol\newline\textbf{?} \emph{Active and second-hand smoking together.} & 872\\
Lead & Lead from paint, dust, water or soil, usually measured as blood lead concentration in children & 249\\
Mercury, arsenic, cadmium and other metals & Toxic metals and metalloids other than lead (mercury, arsenic, cadmium, manganese), measured in blood, urine, hair or water & 464\\
Phthalates, phenols and other endocrine disruptors & Non-persistent endocrine-disrupting chemicals: phthalates, bisphenol A and other phenols, parabens, triclosan & 416\\
PFAS and persistent organic pollutants & Persistent chemicals: per- and polyfluoroalkyl substances, PCBs, organochlorine pesticides such as DDT/DDE, dioxins, flame retardants & 568\\
Pesticides & Current-use pesticides: organophosphates, pyrethroids, herbicides such as glyphosate, and agricultural pesticide application near the home & 931\\
Drinking water contamination & Contaminants in drinking water: nitrate, disinfection by-products, lead service lines, fluoride and water system violations & 209\\
Ambient temperature and heat & Outdoor temperature, heat waves, cold spells and humidity, as a short-term or long-term exposure & 422\\
Environmental noise & Road, rail and aircraft noise at the residence or school & 137\\
Greenness and blue space & Vegetation and water near the home: NDVI greenness, parks, tree canopy, access to blue space & 302\\
Built environment & The physical form of the neighbourhood: walkability, land use, building density, food outlets, urbanicity & 1,259\\
Neighbourhood deprivation and opportunity & Area-level socioeconomic disadvantage measured by a composite index (Area Deprivation Index, Child Opportunity Index, Social Vulnerability Index, census-tract poverty)\newline\textbf{?} \emph{Area-level SDOH is separate from household SDOH (next row). Keep split?} & 2,446\\
Household social and economic factors & Individual or household social determinants: income, insurance type, parental education, food insecurity, housing instability, homelessness & 3,840\\
Structural racism and segregation & Residential segregation, historical redlining, discrimination and other structural measures of racism & 120\\
Violence and crime exposure & Neighbourhood violent crime, gun violence, community violence and adverse childhood experiences & 75\\
Vaccination & Receipt of a vaccine, its type, dose or timing, studied as an exposure that may precede an adverse event or another health...\newline\textbf{?} \emph{One facet for all vaccines. Split by vaccine type?} & 1,324\\
Medication use in pregnancy or early life & A medicine taken in pregnancy or infancy studied as an exposure: antibiotics, acid suppressants, acetaminophen, antidepressants\newline\textbf{?} \emph{Only medicines in pregnancy or infancy. Adult drug exposures excluded.} & 1,681\\
\bottomrule\end{longtable}
\subsection{Health outcome}
A disease, diagnosis, measurement or event studied as the consequence of an exposure.

\begin{longtable}{@{}P{3.6cm}P{9.4cm}r@{}}\toprule Term & Definition (first sentence) & Papers\\\midrule\endhead
Asthma and wheeze & Asthma incidence, prevalence, exacerbation, emergency visit or hospitalisation, and recurrent or persistent wheeze & 766\\
Lung function & Spirometric or other lung function measurements: FEV1, FVC, FEV1/FVC. & 169\\
Allergic disease & Atopic dermatitis, food allergy, allergic rhinitis and allergic sensitisation & 219\\
Obesity and adiposity & BMI, BMI z-score, overweight, obesity, waist circumference or body fat, including rapid infant weight gain & 998\\
Cardiometabolic risk & Blood pressure, hypertension, insulin resistance, type 2 diabetes, dyslipidaemia and metabolic syndrome\newline\textbf{?} \emph{Blood pressure, type 2 diabetes and lipids in one facet.} & 931\\
Type 1 diabetes and autoimmune disease & Type 1 diabetes, coeliac disease, juvenile idiopathic arthritis, inflammatory bowel disease and other autoimmune conditions & 371\\
Neurodevelopment and cognition & Cognitive function, IQ, developmental delay, language and motor development, school performance & 809\\
ADHD and autism & Attention-deficit/hyperactivity disorder and autism spectrum disorder diagnoses or symptom scores\newline\textbf{?} \emph{ADHD and autism share one facet. Split?} & 343\\
Mental and behavioural health & Depression, anxiety, behavioural problems, self-harm, substance use disorder and psychiatric emergency visits & 1,161\\
Preterm birth and gestational age & Birth before 37 weeks, gestational age at birth, and spontaneous preterm labour & 956\\
Birth weight and fetal growth & Birth weight, low birth weight, small or large for gestational age, and fetal growth measured by ultrasound & 919\\
Congenital anomalies & Birth defects: congenital heart defects, neural tube defects, orofacial clefts and others & 516\\
Pregnancy complications & Maternal outcomes of pregnancy: hypertensive disorders including preeclampsia, gestational diabetes, stillbirth and severe maternal morbidity & 419\\
Childhood cancer & Cancer in children and adolescents, chiefly leukaemia, brain tumours and lymphoma\newline\textbf{?} \emph{Only childhood cancer. Adult cancers are tagged nowhere.} & 1,083\\
Infection and infection-related care & Infections and their consequences: respiratory infections, COVID-19, bronchiolitis, otitis media, infection-related hospitalisation & 1,973\\
Vaccine adverse event & An event studied as a possible adverse reaction to a vaccine: febrile seizure, fever, anaphylaxis, Guillain-Barre syndrome, intussusception, myocarditis & 466\\
Mortality & Death from any cause or a specific cause, including infant mortality and sudden infant death\newline\textbf{?} \emph{All-cause and cause-specific mortality together.} & 2,748\\
Health care use & Emergency department visits, hospitalisation, primary care visits or costs, studied as an outcome\newline\textbf{?} \emph{Health care use as an outcome, not as a data source.} & 1,778\\
Puberty timing and linear growth & Age at puberty or menarche, precocious puberty, and height or linear growth & 74\\
\bottomrule\end{longtable}
\subsection{Data source}
Where the individual-level data came from.

\begin{longtable}{@{}P{3.6cm}P{9.4cm}r@{}}\toprule Term & Definition (first sentence) & Papers\\\midrule\endhead
Electronic health records & Electronic health or medical records from clinical care: diagnoses, encounters, vital signs, laboratory results, medications & 2,348\\
Insurance claims & Billing or insurance claims data: Medicaid, Medicare, commercial claims, national health insurance databases & 964\\
Health registry & A population or disease registry: birth, death, cancer, prescription, patient or national health registers & 6,040\\
Immunization registry or vaccine surveillance system & An immunization information system or vaccine safety surveillance system: Vaccine Safety Datalink, VAERS, provincial vaccine registries & 440\\
Record linkage & Individual-level linkage of two or more data sources, such as hospital records to a birth registry or to environmental data by address & 443\\
Research cohort & Data collected for research by a prospective cohort, most often a birth or mother-child cohort with its own visits, samples and questionnaires & 2,914\\
National health survey & A repeated cross-sectional national survey with examination or biomarker data, such as NHANES. & 305\\
Biobank & A biobank linking stored samples or genotypes to health records, such as UK Biobank or All of Us & 279\\
\bottomrule\end{longtable}
\subsection{Exposure assessment method}
How an exposure was assigned to a person: from an address, a model, a monitor, a biological sample, a questionnaire or a record.

\begin{longtable}{@{}P{3.6cm}P{9.4cm}r@{}}\toprule Term & Definition (first sentence) & Papers\\\midrule\endhead
Residential address geocoding & Assigning area-level exposures from a geocoded home address, census tract or postcode & 1,332\\
Spatial or satellite exposure model & An exposure surface estimated by land-use regression, dispersion, chemical transport, satellite remote sensing or machine-learning models & 311\\
Fixed monitoring station & Exposure taken from the nearest regulatory air quality or weather monitor & 583\\
Biomonitoring & Measuring a chemical or its metabolite in blood, urine, hair, breast milk or another biological sample & 1,324\\
Questionnaire or interview & Exposure reported by the participant or a parent in a questionnaire or interview & 2,062\\
Personal sensor or wearable & A personal monitor, wearable sensor, smartphone or GPS tracking used to measure an individual's exposure directly & 243\\
Exposure taken from health records & An exposure read from the health record itself: a vaccination record, a prescription, a recorded smoking status or a screening test result & 298\\
\bottomrule\end{longtable}
\subsection{Study design or analysis}
The design or analytic approach.

\begin{longtable}{@{}P{3.6cm}P{9.4cm}r@{}}\toprule Term & Definition (first sentence) & Papers\\\midrule\endhead
Exposome-wide association design & Testing many exposures agnostically against one outcome with correction for multiple testing (EWAS, ExWAS), including polyexposure risk scores & 948\\
Exposure mixture methods & Statistical methods for joint effects of correlated exposures: weighted quantile sum, Bayesian kernel machine regression, quantile g-computation & 229\\
Self-controlled design & A design in which each person is their own control: self-controlled case series, case-crossover, risk-interval designs & 415\\
Time-series design & Daily or weekly counts of events regressed on short-term exposure variation, often with distributed lags & 331\\
Machine learning & Random forests, gradient boosting, neural networks, LASSO or other machine learning used for prediction or variable selection & 554\\
Gene-environment interaction & Joint analysis of genotypes and exposures, including GxE interaction tests and polygenic-by-exposure models & 300\\
\bottomrule\end{longtable}
\subsection{Age stratum}
The life stage the study enrolled or the exposure window it studied.

\begin{longtable}{@{}P{3.6cm}P{9.4cm}r@{}}\toprule Term & Definition (first sentence) & Papers\\\midrule\endhead
Pregnancy or prenatal window & Exposure in pregnancy, or pregnant people enrolled & 3,697\\
Infants & Newborns and infants under two years old & 3,236\\
Children & Children from two to twelve years old & 4,356\\
Adolescents & Young people from thirteen to eighteen years old & 2,629\\
Adults & Adults, including older adults & 7,290\\
\bottomrule\end{longtable}
\paragraph{Not in the vocabulary.} Diet and nutrition, physical activity, sleep, the microbiome, ionising
radiation, occupational exposures, adult cancers, and kidney or liver outcomes. Adding any of them is a new facet in
\texttt{config/facets.yaml}; papers are re-tagged without re-searching.
\decision{Approve the term lists? Mark terms to split, merge, rename or add. The flagged rows above are the
ones I was least sure of.}
\section{Decision 4: per-paper questions}\label{sec:coding}
These replace the old free-text checklist (\texttt{config/study\_coding.yaml}). Each is a yes/no or a fixed set of
options, answered from the passages listed, with the probability and the paragraph that carried the answer kept.
Summary, key findings and limitations are no longer produced: they are not decisions.

\begin{center}\begin{longtable}{@{}lP{1.4cm}P{2.4cm}P{8.2cm}@{}}\toprule
Field & Type & Reads & Options\\\midrule\endhead
\texttt{primary\_research} & yes no & abstract & yes / no\\
\texttt{routinely\_collected\_data} & yes no & abstract, methods & yes / no\\
\texttt{data\_sources} & any of & abstract, methods & the data source terms of Section~\ref{sec:vocab}\\
\texttt{age\_strata} & any of & abstract, methods & the population terms of Section~\ref{sec:vocab}\\
\texttt{study\_design} & one of & abstract, methods & cohort; case\_control; cross\_sectional; self\_controlled; time\_series; trial; exposome\_wide\\
\texttt{exposures\_studied} & any of & abstract & the exposure terms of Section~\ref{sec:vocab}\\
\texttt{outcomes\_studied} & any of & abstract & the outcome terms of Section~\ref{sec:vocab}\\
\texttt{data\_availability} & one of & availability, methods & public\_repository; on\_request; restricted; supplementary\_only\\
\bottomrule\end{longtable}\end{center}
\texttt{routinely\_collected\_data} keeps the old \texttt{ehr\_used} definition: electronic, individual-level and
collected for care, billing or administration before the study. Chart review by study staff, research questionnaires
and aggregate counts do not count. The pediatric EHR corpus is then: primary research = yes, routinely collected
data = yes, and age strata include infant, child or adolescent.
\decision{Are these the fields your review needs? Anything to add (e.g.\ sample size band, country, outcome
definition by codes vs.\ clinical measure, exposure window)? Is the \texttt{ehr\_used} definition still right?}
\section{What does not need your review}\label{sec:skip}
\begin{itemize}
\item The harvester, parser, store and concept resolvers: ported from POTS-phenotyping with their tests.
\item The decision-model and Claude-judge code: tested offline with a fake model; also run end to end over the 255
seed paragraphs.
\item Derived tables (\texttt{data/derived/}) and the harvest report: generated; 24 of 26 checks pass. The two failures
are 19 terms with no OMOP concept match, mostly exposure concepts that the free SNOMED lookup does not cover; an
ATHENA vocabulary download fixes most of them.
\item The move of the old pipeline to \texttt{legacy/}.
\end{itemize}

\section*{After your review}
\begin{enumerate}
\item Apply your changes to the config files; re-harvest only if Decision 1 changed.
\item Run the seed reader run on the lab server and the Claude judge (\texttt{docs/DECISION\_MODELS.md}).
\item Hand-code 30--50 papers for the Decision 4 fields to measure agreement before the coding is used.
\item Before merging: the GitHub Pages site (old site moved to \texttt{legacy/docs/}), and the POTS branch this is
built on, which is not yet merged.
\end{enumerate}

\appendix
\section{Link predicates}
\begin{center}\begin{tabular}{@{}lP{4.2cm}P{7.8cm}@{}}\toprule Predicate & From $\to$ to & Question option\\\midrule
\texttt{risk\_factor\_for} & exposure $\to$ outcome & an exposure associated with more of, or a higher risk of, an outcome\\
\texttt{protective\_for} & exposure $\to$ outcome & an exposure associated with less of, or a lower risk of, an outcome\\
\texttt{no\_association\_with} & exposure $\to$ outcome & an exposure found not to be associated with an outcome\\
\texttt{assessed\_by} & exposure $\to$ assessment & how an exposure was measured or assigned\\
\texttt{ascertained\_from} & outcome, exposure $\to$ data\_source & the data source an outcome or exposure was taken from\\
\texttt{linked\_with} & data\_source, assessment $\to$ data\_source, assessment & two data sources joined person by person, or by address\\
\bottomrule\end{tabular}\end{center}
\end{document}
